\section{The two optimizers}
\label{s-optimizers}

The relaxed and integer problems both reduce to one-dimensional
constructions. A scalar equation gives the relaxed optimizer. A
fractional-knapsack ordering gives every integer optimizer and a finite
algorithm for computing it.

\subsection{The relaxed optimizer}
\label{s-relaxed-scalar}

For \(T>0\), define the log-size of the coordinate response by
\[
    s_m(T)
    =
    \sum_{i=1}^{m}
    \max\{\log p_i,\log(1+T)-\log p_i\}.
\]

\begin{proposition}[Relaxed optimizer]
\label{prop-scalar-reduction}
The relaxed problem has a unique maximizer \(\rstar\). There is a unique
solution \(T^\star>0\) of
\[
    T=s_m(T)\log s_m(T),
\]
and
\begin{equation}
\label{e-relaxed-response}
    r_i^\star
    =
    \max\left\{
        1,\,
        \frac{\log(1+T^\star)}{\log p_i}-1
    \right\}.
\end{equation}
In particular,
\[
    S^\star:=S(\rstar)=s_m(T^\star),
    \qquad
    T^\star=S^\star\log S^\star.
\]
\end{proposition}

\begin{proof}
The local logarithmic terms in \eqref{e-objective} are bounded above,
whereas \(\log\log S(r)\) tends to infinity whenever a coordinate of
\(r\) does. A maximum therefore exists. At any maximum, put
\(T=S(r)\log S(r)\). Since
\[
    \frac{d}{du}\log(p_i-p_i^{-u})
    =
    \frac{\log p_i}{p_i^{u+1}-1},
\]
the one-sided first-order conditions give
\[
    p_i^{r_i+1}-1=T
    \quad\hbox{if \(r_i>1\)},\qquad
    p_i^2-1\geq T
    \quad\hbox{if \(r_i=1\)}.
\]
These conditions are equivalent to \eqref{e-relaxed-response}, and the
size of that response is \(s_m(T)\).

It remains to show that the scalar equation has one solution. Set
\[
    h(T)=T-s_m(T)\log s_m(T).
\]
We have \(h(0)=-\theta(p_m)\log\theta(p_m)<0\), while
\(s_m(T)\leq\theta(p_m)+m\log(1+T)\), so \(h(T)\) tends to infinity.
Every root satisfies \(T\geq\theta(p_m)\log\theta(p_m)\).

Away from the switching points, let \(k\) be the number of primes with
\(p_i^2<1+T\). If \(k=0\), then \(h'(T)=1\). Otherwise,
\[
    h'(T)
    =
    1-\frac{k(1+\log s_m(T))}{1+T}.
\]
The \(k\)th prime is at least \(k+1\), so \(1+T>(k+1)^2\). Also
\[
    s_m(T)
    =
    \theta(p_m)
    +
    \sum_{p_i^2<1+T}
    \bigl(\log(1+T)-2\log p_i\bigr)
    <2(1+T).
\]
Indeed, on the root interval \(\theta(p_m)<1+T\), while
\(k<\sqrt{1+T}\) and \(\log(1+T)<\sqrt{1+T}\).
The elementary inequality
\[
    k\bigl(1+\log(2u)\bigr)<u,
    \qquad u>(k+1)^2,
\]
now gives \(h'(T)>0\). At the left endpoint the inequality reduces to
\[
    k+1+\frac1k-\log2-2\log(k+1)>0.
\]
This is immediate for \(k=1,2\) and increasing for \(k\geq2\).
Continuity across the switching points proves uniqueness of the root.
Every global maximizer satisfies the first-order conditions, so the
response at this root is the unique maximizer.
\end{proof}

\subsection{The integer optimizer}
\label{s-integer-prefix}

For \(x>1\) and an integer \(k\geq1\), define the standard CA score
\cite[Section~4.1]{nicolas-2022-sigma}
\begin{equation}
\label{e-ca-score}
    F(x,k)
    =
    \frac1{\log x}
    \log\left(1+\frac1{x+x^2+\cdots+x^k}\right).
\end{equation}
Raising the exponent of \(p\) from \(k-1\) to \(k\) changes the log
divisor ratio by
\begin{equation}
\label{e-score-dictionary}
    \log\left(
        \frac{\sigma(p^k)/p^k}
             {\sigma(p^{k-1})/p^{k-1}}
    \right)
    =
    \log p\,F(p,k).
\end{equation}
Equivalently,
\begin{equation}
\label{e-score-integral}
    F(p,k)
    =
    \int_{k-1}^{k}\frac{du}{p^{u+1}-1}.
\end{equation}
This shows that the score decreases strictly in both arguments and that
\[
    F(p,k)\leq\frac1{p^k-1}.
\]
In particular, only finitely many eligible scores exceed any positive
level.
The fixed-support problem includes the increments \((p,k)\) with
\(p\leq p_m\) and \(k\geq2\). Their log-size weights are \(\log p\),
and their scores are ordered as above.

Group equal scores and order the resulting blocks by decreasing score.
Starting from the primorial \(\prod_{p\leq p_m}p\), let \(N_j\) be the
integer obtained after the first \(j\) blocks. We refer to these
integers as score-ordered prefixes. Every such prefix
is a valid exponent vector because each predecessor in a prime chain has
a larger score.

\begin{theorem}[Prefix characterization]
\label{thm-integer-prefix}
Every integer maximizer is one of the prefix integers \(N_j\). If \(N_m^\star\)
is any maximizer and
\[
    L_m=\log N_m^\star,
    \qquad
    \eps_m=\frac1{L_m\log L_m},
\]
then an increment \((p,k)\), \(k\geq2\), is selected exactly when
\begin{equation}
\label{e-integer-threshold}
    F(p,k)>\eps_m.
\end{equation}
No eligible score equals \(\eps_m\), and no maximizer cuts a tie block.
Every maximizing prefix also satisfies
\begin{equation}
\label{e-prefix-bound}
    \log N_m^\star\leq\theta(p_m)^{\zeta(2)}.
\end{equation}
\end{theorem}

\begin{proof}
Consider the piecewise-linear curve through the points
\[
    \left(
        \log N_j,\,
        \log\frac{\sigma(N_j)}{N_j}
    \right).
\]
Its slopes are the ordered scores. This curve is a fractional-knapsack
upper bound. Let \(\mathcal E\) be any finite selection of increments
with total extra log-size \(z\). On an interval with score \(c\), let
\(\mathcal P\) be the set of all increments in earlier blocks. Then
\[
\begin{aligned}
    \sum_{(p,k)\in\mathcal E}\log p\,F(p,k)
    &=
    cz+
    \sum_{(p,k)\in\mathcal E}
    \log p\bigl(F(p,k)-c\bigr)\\
    &\leq
    cz+
    \sum_{(p,k)\in\mathcal P}
    \log p\bigl(F(p,k)-c\bigr).
\end{aligned}
\]
The right-hand side is the linear interpolation between the two adjacent
score-ordered prefix values.

On each interpolation interval the log divisor ratio is affine in
\(\log n\), while \(-\log\log\log n\) is strictly convex. One of the two
prefix endpoints therefore has at least as large a value of \(\log G\)
as any integer selection in the interval. Thus the supremum over integer
selections equals the supremum over score-ordered prefixes.

The total possible gain above the primorial is
\[
    -\sum_{p\leq p_m}\log(1-p^{-2})
    \leq
    \log\zeta(2).
\]
It follows that the prefix objective tends to minus infinity, so a
maximizing prefix exists. Comparing any maximizing prefix with the
primorial gives
\[
    \log\log\log N_m^\star
    \leq
    \log\log\theta(p_m)+\log\zeta(2),
\]
which is \eqref{e-prefix-bound}.

It remains to characterize an arbitrary maximizer \(N_m^\star\). The function
\(u\mapsto\log\log u\) is strictly concave, and its derivative at \(L_m\)
is \(\eps_m\). Optimality and the tangent inequality give, for every
integer \(M\) on the prescribed support,
\[
\begin{aligned}
    \log\frac{\sigma(M)}M-\log\log\log M
    &\leq
    \log\frac{\sigma(N_m^\star)}{N_m^\star}-\log\log L_m,\\
    \log\log\log M
    &\leq
    \log\log L_m+\eps_m(\log M-L_m).
\end{aligned}
\]
Consequently,
\[
    \log\frac{\sigma(M)}M-\eps_m\log M
    \leq
    \log\frac{\sigma(N_m^\star)}{N_m^\star}-\eps_m L_m.
\]
This expression is separable over the optional increments. Equation
\eqref{e-score-dictionary} shows that an increment is selected when its
score exceeds \(\eps_m\) and is unselected when its score is smaller.
Equality cannot occur. Adding or removing an increment with score
\(\eps_m\) would preserve the last display, while strictness in the
tangent inequality would produce a larger value of \(\log G\). Thus the
selected increments form a score-ordered prefix and satisfy
\eqref{e-integer-threshold}. This also proves the tie assertion and
extends \eqref{e-prefix-bound} to every maximizer.
\end{proof}

With exact score comparisons, the theorem gives a finite algorithm.
Maintain one score-ordered chain for each prime, repeatedly take
the largest available block, and evaluate \(G\) at the resulting prefix.
The search stops once \(\log N_j\) exceeds the bound in
\eqref{e-prefix-bound}. This computes the optimizer and its strict
threshold without searching an exponent box.

The threshold also gives the precise SA and CA interpretation. The
integer \(N_m^\star\) maximizes
\[
    \frac{\sigma(n)}{n^{1+\eps_m}}
\]
over integers on the prescribed support. It consequently maximizes
\(\sigma(n)/n\) on the same support under the size constraint
\(\log n\leq L_m\). It need not be a global CA or SA number because
the first-exponent increments \(F(p,1)\) are fixed by the support
constraint.

\subsection{The common scale}
\label{s-common-scale}

\begin{proposition}[Relaxed scale]
\label{prop-relaxed-scale}
As \(m\) tends to infinity,
\[
    S^\star\sim\theta(p_m)\sim p_m,
    \qquad
    T^\star\sim p_m\log p_m,
\]
and
\begin{equation}
\label{e-relaxed-displacement}
    S^\star-\theta(p_m)
    =
    O\left(\sqrt{\frac{p_m}{\log p_m}}\right).
\end{equation}
\end{proposition}

\begin{proof}
Put \(y=\sqrt{1+T^\star}\). We first show that \(y\leq p_m\) for all
sufficiently large \(m\). Otherwise every coordinate is interior, so
\[
    S^\star
    =
    2m\log y-\theta(p_m)
    \leq2m\log y.
\]
Since \(y^2-1=S^\star\log S^\star<2S^\star\log y\), this would imply
\(y/\log y\ll\sqrt m\). But \(y>p_m\) and the prime number theorem give
\(y/\log y>p_m/\log p_m\sim m\), a contradiction.

Splitting the response \eqref{e-relaxed-response} at \(p=y\) gives
\[
\begin{aligned}
    S^\star-\theta(p_m)
    &=
    2\bigl(\pi(y)\log y-\theta(y)\bigr)\\
    &=
    2\int_2^y\frac{\pi(u)}u\,du
    =
    O\left(\frac{y}{\log y}\right).
\end{aligned}
\]
Because \(y^2-1=S^\star\log S^\star\), the last bound is \(o(S^\star)\).
It follows that \(S^\star\sim\theta(p_m)\sim p_m\), and hence
\(T^\star\sim p_m\log p_m\). Substitution in the same bound proves
\eqref{e-relaxed-displacement}.
\end{proof}

\begin{proposition}[Integer scale]
\label{prop-integer-threshold-scale}
For every integer maximizer \(N_m^\star\),
\begin{equation}
\label{e-integer-mass}
    \log N_m^\star-\theta(p_m)
    =
    \bigl(\sqrt2+o(1)\bigr)
    \sqrt{\frac{T^\star}{\log T^\star}},
\end{equation}
and
\begin{equation}
\label{e-threshold-separation}
    1-T^\star\eps_m
    =
    \bigl(\sqrt2+o(1)\bigr)
    \sqrt{\frac{\log T^\star}{T^\star}}.
\end{equation}
In particular, \(T^\star\eps_m<1\) for all sufficiently large \(m\).
\end{proposition}

\begin{proof}
Put \(L_m=\log N_m^\star\). Strict separation in
Theorem~\ref{thm-integer-prefix} gives the exact identity
\[
    L_m-\theta(p_m)
    =
    \sum_{p\leq p_m}\log p\,
    \#\{k\geq2\mid F(p,k)>\eps_m\}.
\]
The bound \(F(p,k)\leq1/(p^k-1)\) first gives
\[
    L_m-\theta(p_m)
    =
    O\left(\sqrt{L_m\log L_m}\right).
\]
Since \(L_m=\theta(p_m)+O(\sqrt{L_m\log L_m})\), we have
\(L_m\sim\theta(p_m)\sim p_m\). Hence
\[
    L_m\log L_m\sim T^\star.
\]
We now estimate the score mass at this level. Put
\[
    U=L_m\log L_m.
\]
The first optional increment satisfies
\[
    F(p,2)
    =
    \frac{1+o(1)}{p^2\log p}.
\]
For every fixed \(\delta>0\) and every sufficiently large \(p\), this
gives the uniform bounds
\[
    \frac{1-\delta}{p^2\log p}
    \leq
    F(p,2)
    \leq
    \frac{1+\delta}{p^2\log p}.
\]
The positive solution of \(x^2\log x=U\) is
\[
    \bigl(\sqrt2+o(1)\bigr)
    \sqrt{\frac{U}{\log U}}.
\]
Comparison with the positive solutions of
\(x^2\log x=(1\pm\delta)U\), followed by the prime number theorem and
then \(\delta\to0\), gives
\[
    \sum_{\substack{p\leq p_m\\ F(p,2)>1/U}}\log p
    =
    \bigl(\sqrt2+o(1)\bigr)
    \sqrt{\frac{U}{\log U}}.
\]
The support cutoff is inactive because
\[
    p_m\asymp\frac{T^\star}{\log T^\star}
    \gg\sqrt{\frac{U}{\log U}}.
\]
For \(k\geq3\), the bound \(F(p,k)\leq1/(p^k-1)\) gives total log weight
\[
    \sum_{3\leq k\leq\log_2(1+U)}
    \theta\left((1+U)^{1/k}\right)
    =
    O(U^{1/3}+U^{1/4}\log U),
\]
which is lower order. Since \(U\sim T^\star\), the exact score identity
now proves \eqref{e-integer-mass}.

By \eqref{e-relaxed-displacement},
\[
    S^\star-\theta(p_m)
    =
    o\left(
        \sqrt{\frac{T^\star}{\log T^\star}}
    \right).
\]
Together with \eqref{e-integer-mass}, this gives
\[
    L_m-S^\star
    =
    \bigl(\sqrt2+o(1)\bigr)
    \sqrt{\frac{T^\star}{\log T^\star}}.
\]
Taylor expansion of \(u\log u\) now gives
\[
\begin{aligned}
    L_m\log L_m-T^\star
    &=
    (1+\log S^\star)(L_m-S^\star)
    +O\left(\frac{(L_m-S^\star)^2}{S^\star}\right)\\
    &=
    \bigl(\sqrt2+o(1)\bigr)
    \sqrt{T^\star\log T^\star}.
\end{aligned}
\]
Taking reciprocals proves \eqref{e-threshold-separation}.
\end{proof}
